\subsection{非隣接2項型}\label{節：非隣接2項型}

非隣接2項間\bookterm{recurrence}{漸化式}
\[
  a_{n+2}=f(a_n)\qquad(n\ge1)
\]
では,$a_1$から奇数番目,$a_2$から偶数番目が決まる.
$u_k=a_{2k-1},v_k=a_{2k}$と部分列に分ける.
\[
  \begin{array}{c|cccc}
    k&1&2&3&\cdots\\ \hline
    u_k&a_1&a_3&a_5&\cdots\\
    v_k&a_2&a_4&a_6&\cdots
  \end{array}
\]
$n$が2進むと$k$は1進み,$u_{k+1}=f(u_k),v_{k+1}=f(v_k)$となる.
最後に奇数$n$では$k=\dfrac{n+1}{2}$,偶数$n$では$k=\dfrac{n}{2}$で戻す.

また,$a_{n+2}=qa_n$では,$a_{n+2}=0\cdot a_{n+1}+qa_n$と見て\bookterm{characteristic}{特性方程式}を使うこともできる.

\noindent \bookblocklink{example-nonadjacent-1}{problem}{answer}{【例題】}\par\nobreak
\begin{enumerate}[label=(\arabic*),parsep=3pt,format=\bookitemlink{example-nonadjacent-1}{problem}{answer}]
  \item $\displaystyle a_1=2,\,a_2=3,\,a_{n+2}=a_n^2$
  \item $\displaystyle a_1=1,\,a_2=2,\,a_{n+2}=2a_n+1$
  \item $\displaystyle a_1=1,\,a_2=3,\,a_{n+2}=4a_n$
\end{enumerate}

\noindent\bookblocklink{example-nonadjacent-1}{hint}{problem}{【方針】}\par\nobreak
\begin{enumerate}[label=(\arabic*),parsep=3pt,format=\bookitemlink{example-nonadjacent-1}{hint}{problem}]
  \item 奇数番目と偶数番目を別々の\bookterm{sequence}{数列}とみなし,それぞれの二乗関係を調べます.
  \item $b_n=a_n+1$とおいて定数項を消してから,奇数番目と偶数番目に分けます.
  \item $a_{n+2}=0\cdot a_{n+1}+4a_n$と見て,\bookterm{characteristic}{特性方程式}を解きます.
\end{enumerate}

\noindent\bookblocklink{example-nonadjacent-1}{answer}{problem}{【解答】}\par\nobreak
\begin{enumerate}[label=(\arabic*),parsep=3pt,format=\bookitemlink{example-nonadjacent-1}{answer}{problem}]
  \item $\displaystyle a_{2k-1}=2^{\,2^{k-1}},\quad a_{2k}=3^{\,2^{k-1}}\quad(k\ge1)$
  \item $\displaystyle a_{2k-1}=2^k-1,\quad a_{2k}=3\cdot2^{k-1}-1\quad(k\ge1)$
  \item $\displaystyle a_n=\frac54\,2^{n-1}-\frac14\,(-2)^{n-1}$
\end{enumerate}

\noindent\bookblocklink{example-nonadjacent-1}{solution}{problem}{【解法】}\par\nobreak
\begin{enumerate}[label=(\arabic*),parsep=3pt,format=\bookitemlink{example-nonadjacent-1}{solution}{problem}]
  \item $u_k=a_{2k-1},\,v_k=a_{2k}$とおくと,
    \[
      u_1=2,\quad v_1=3,\qquad
      u_{k+1}=u_k^2,\quad v_{k+1}=v_k^2
    \]
    となる.したがって,それぞれ順に2乗していくので,
    \[
      u_k=2^{\,2^{k-1}},\qquad v_k=3^{\,2^{k-1}}
    \]
    となり,$u_k=a_{2k-1},v_k=a_{2k}$で戻す.
  \item $b_n=a_n+1$とおくと,
    \[
      b_{n+2}=a_{n+2}+1=2a_n+2=2b_n
    \]
    であり,$b_1=2,\,b_2=3$である.
    よって,奇数番目と偶数番目で分けると,
    \[
      b_{2k-1}=2^k,\qquad b_{2k}=3\cdot2^{k-1}
    \]
    となる.$a_n=b_n-1$に戻せば【解答】を得る.
  \item \bookterm{characteristic}{特性方程式}は
    \[
      r^2-4=0
    \]
    であり,解は$r=2,-2$である.したがって,
    \[
      a_n=A\,2^{n-1}+B\,(-2)^{n-1}
    \]
    とおける.$a_1=1,\,a_2=3$を代入すると,
    \[
      A+B=1,\qquad 2A-2B=3
    \]
    より$A=\dfrac{5}{4},B=-\dfrac{1}{4}$と決まり,【解答】を得る.
\end{enumerate}

\subsubsection{一周期分をまとめる}\label{節：一周期分をまとめる}
更新の係数が奇数回と偶数回で変わる場合も,二回分をまとめられる.

\bookstep{例題：二回分をまとめる}
\[
  a_1=1,\qquad
  a_{n+1}=\begin{cases}
    2a_n+1&(n\text{が奇数}),\\
    3a_n+2&(n\text{が偶数})
  \end{cases}
\]
奇数番目を$x_k=a_{2k-1}$と置くと,$a_{2k}=2x_k+1$であり,
次の更新まで続けて
\[
  x_{k+1}=3(2x_k+1)+2=6x_k+5
\]
を得る.従って$x_1=1$から$x_k=2\cdot6^{k-1}-1$であり,
\BookMathLayout{\[
    a_{2k-1}=2\cdot6^{k-1}-1,\qquad
    a_{2k}=4\cdot6^{k-1}-1\quad(k\ge1)
  \]}{\[
    \begin{aligned}
      a_{2k-1}&=2\cdot6^{k-1}-1,\\
      a_{2k}&=4\cdot6^{k-1}-1\quad(k\ge1)
    \end{aligned}
  \]}
となる.偶数番目は更新の途中から戻した.
付加項も次の倍率を受けるため,二回分は$1+2$でなく$3\cdot1+2=5$である.

\bookstep{一周期への拡張}
$a_{n+1}=p_na_n+q_n$の係数が周期$L$をもつなら,
$x_k=a_{1+kL}$ $(k\ge0)$と置き,$x_0=a_1$から一周期分を合成する.
\[
  x_{k+1}=Px_k+Q
\]
の$P$は倍率の積,$Q$は各付加項にその後の倍率を掛けた寄与の和である.\footnote{最初の一周期の係数を$p_1,\ldots,p_L$, $q_1,\ldots,q_L$とすると,
  $P=\prod_{j=1}^L p_j$, $Q=\sum_{j=1}^L q_j\prod_{i=j+1}^L p_i$である.空積は$1$とする.
  この表示は倍率に$0$を含む場合も成り立つ.}
途中の項は元の一回ずつの更新から戻せる.
$P=0$なら$x_k=Q$ $(k\ge1)$,$P=1$なら$x_k=x_0+kQ$である.
それ以外は,\sectionref*{節：特性方程式型}と同じく不動点からの差を取る.

\begin{bookpoint}
  \textbf{更新の間隔を選ぶ}\par
  係数の一周期分をまとめる.係数の周期性は,項の値の周期性とは異なる.
\end{bookpoint}
付録の\sectionref*{節：二回分の更新}では,同じ計算を行列の積で読む.
符号も変わる例は\bookproblemref{practice-connections}{4}{接続演習問4}へ.
